\begin{proofbox}

	\textbf{\textcolor{CProof}{思路提示}}

	旋转曲面的面积公式为
	\[
		S
		=
		2\pi\int y\sqrt{1+(y')^{2}}\,dx.
	\]

	将椭圆方程
	\[
		\frac{x^{2}}{a^{2}}+\frac{y^{2}}{b^{2}}=1
	\]
	代入并化简，再根据 $a$ 与 $b$ 的大小关系分类讨论。

	当 $a>b$ 时，用三角代换，得到含 $\arcsin$ 的公式；当 $a<b$ 时，用双曲代换，得到含 $\operatorname{artanh}$ 的公式；当 $a=b$ 时，退化为球，直接得到 $4\pi R^{2}$。

	\medskip

	\textbf{\textcolor{CProof}{正式推导}}

	\begin{description}[style=nextline, leftmargin=2.8em, labelsep=0.5em, itemsep=0.55em]

		\item[\textbf{步骤一（建立积分）：}]
		      椭圆上半部分可以写成
		      \[
			      y
			      =
			      b\sqrt{1-\frac{x^{2}}{a^{2}}}.
		      \]

		      它绕 $x$ 轴旋转一周形成旋转椭球面。

		      根据旋转曲面面积公式，有
		      \[
			      S
			      =
			      2\pi\int_{-a}^{a}
			      y\sqrt{1+(y')^{2}}\,dx.
		      \]

		\item[\textbf{步骤二（化简被积函数）：}]
		      由
		      \[
			      y
			      =
			      b\sqrt{1-\frac{x^{2}}{a^{2}}}
			      =
			      \frac{b}{a}\sqrt{a^{2}-x^{2}},
		      \]
		      可计算出 $y'$，并进一步化简得到
		      \[
			      y\sqrt{1+(y')^{2}}
			      =
			      \frac{b}{a^{2}}
			      \sqrt{a^{4}-(a^{2}-b^{2})x^{2}}.
		      \]

		      因此
		      \[
			      S
			      =
			      \frac{2\pi b}{a^{2}}
			      \int_{-a}^{a}
			      \sqrt{a^{4}-(a^{2}-b^{2})x^{2}}\,dx.
		      \]

		      记
		      \[
			      I
			      =
			      \int_{-a}^{a}
			      \sqrt{a^{4}-(a^{2}-b^{2})x^{2}}\,dx,
		      \]
		      则
		      \[
			      S
			      =
			      \frac{2\pi b}{a^{2}}I.
		      \]

		\item[\textbf{步骤三（分类讨论）：}]
		      令
		      \[
			      \Delta=a^{2}-b^{2}.
		      \]

		      根据 $\Delta$ 的符号，分三种情况：
		      \[
			      \Delta>0,\qquad \Delta<0,\qquad \Delta=0.
		      \]

	\end{description}

	\medskip

	\textbf{\textcolor{CProof}{情形一：$a>b$}}

	\begin{description}[style=nextline, leftmargin=2.8em, labelsep=0.5em, itemsep=0.55em]

		\item[\textbf{步骤四（三角代换）：}]
		      当 $a>b$ 时，令
		      \[
			      c=\sqrt{a^{2}-b^{2}}>0.
		      \]

		      则
		      \[
			      e=\frac{c}{a}
			      =
			      \sqrt{1-\frac{b^{2}}{a^{2}}}.
		      \]

		      此时
		      \[
			      I
			      =
			      \int_{-a}^{a}
			      \sqrt{a^{4}-c^{2}x^{2}}\,dx.
		      \]

		      由于被积函数是偶函数，所以
		      \[
			      I
			      =
			      2\int_{0}^{a}
			      \sqrt{a^{4}-c^{2}x^{2}}\,dx.
		      \]

		      作三角代换
		      \[
			      x=\frac{a^{2}}{c}\sin\theta,
			      \qquad
			      dx=\frac{a^{2}}{c}\cos\theta\,d\theta.
		      \]

		      当 $x=0$ 时，$\theta=0$；当 $x=a$ 时，
		      \[
			      \sin\theta=\frac{c}{a}=e,
			      \qquad
			      \theta=\arcsin e.
		      \]

		      同时
		      \[
			      \sqrt{a^{4}-c^{2}x^{2}}
			      =
			      a^{2}\cos\theta.
		      \]

		      因此
		      \[
			      I
			      =
			      \frac{2a^{4}}{c}
			      \int_{0}^{\arcsin e}
			      \cos^{2}\theta\,d\theta.
		      \]

		\item[\textbf{步骤五（积分回代）：}]
		      因为
		      \[
			      \int \cos^{2}\theta\,d\theta
			      =
			      \frac12(\theta+\sin\theta\cos\theta),
		      \]
		      所以
		      \[
			      \int_{0}^{\arcsin e}
			      \cos^{2}\theta\,d\theta
			      =
			      \frac12
			      \left(
			      \arcsin e
			      +
			      e\cdot \frac{b}{a}
			      \right).
		      \]

		      于是
		      \[
			      I
			      =
			      \frac{a^{4}}{c}
			      \left(
			      \arcsin e+\frac{eb}{a}
			      \right).
		      \]

		      代入
		      \[
			      S=\frac{2\pi b}{a^{2}}I
		      \]
		      并利用 $c=ae$，可得
		      \[
			      S
			      =
			      2\pi b^{2}
			      +
			      \frac{2\pi ab}{e}\arcsin e.
		      \]

		      因此，当 $a>b$ 时，
		      \[
			      \boxed{
				      S
				      =
				      2\pi b^{2}
				      +
				      \frac{2\pi ab}{e}\arcsin e
			      },
			      \qquad
			      e=\sqrt{1-\frac{b^{2}}{a^{2}}}.
		      \]

	\end{description}

	\medskip

	\textbf{\textcolor{CProof}{情形二：$a<b$}}

	\begin{description}[style=nextline, leftmargin=2.8em, labelsep=0.5em, itemsep=0.55em]

		\item[\textbf{步骤六（双曲代换）：}]
		      当 $a<b$ 时，令
		      \[
			      d=\sqrt{b^{2}-a^{2}}>0.
		      \]

		      则
		      \[
			      e=\frac{d}{b}
			      =
			      \sqrt{1-\frac{a^{2}}{b^{2}}}.
		      \]

		      此时
		      \[
			      I
			      =
			      \int_{-a}^{a}
			      \sqrt{a^{4}+d^{2}x^{2}}\,dx.
		      \]

		      由于被积函数是偶函数，所以
		      \[
			      I
			      =
			      2\int_{0}^{a}
			      \sqrt{a^{4}+d^{2}x^{2}}\,dx.
		      \]

		      作双曲代换
		      \[
			      x=\frac{a^{2}}{d}\sinh t,
			      \qquad
			      dx=\frac{a^{2}}{d}\cosh t\,dt.
		      \]

		      当 $x=0$ 时，$t=0$；当 $x=a$ 时，
		      \[
			      \sinh t=\frac{d}{a},
			      \qquad
			      t=\operatorname{arsinh}\frac{d}{a}.
		      \]

		      同时
		      \[
			      \sqrt{a^{4}+d^{2}x^{2}}
			      =
			      a^{2}\cosh t.
		      \]

		      因此
		      \[
			      I
			      =
			      \frac{2a^{4}}{d}
			      \int_{0}^{\operatorname{arsinh}(d/a)}
			      \cosh^{2}t\,dt.
		      \]

		\item[\textbf{步骤七（积分回代）：}]
		      因为
		      \[
			      \int \cosh^{2}t\,dt
			      =
			      \frac12(t+\sinh t\cosh t),
		      \]
		      所以
		      \[
			      \int_{0}^{\operatorname{arsinh}(d/a)}
			      \cosh^{2}t\,dt
			      =
			      \frac12
			      \left(
			      \operatorname{arsinh}\frac{d}{a}
			      +
			      \frac{bd}{a^{2}}
			      \right).
		      \]

		      于是
		      \[
			      I
			      =
			      \frac{a^{4}}{d}
			      \left(
			      \operatorname{arsinh}\frac{d}{a}
			      +
			      \frac{bd}{a^{2}}
			      \right).
		      \]

		      代入
		      \[
			      S=\frac{2\pi b}{a^{2}}I
		      \]
		      得到
		      \[
			      S
			      =
			      2\pi b^{2}
			      +
			      \frac{2\pi ba^{2}}{d}
			      \operatorname{arsinh}\frac{d}{a}.
		      \]

		      又因为
		      \[
			      d=be,
			      \qquad
			      \operatorname{arsinh}\frac{d}{a}
			      =
			      \operatorname{artanh}e,
		      \]
		      所以
		      \[
			      S
			      =
			      2\pi b^{2}
			      +
			      \frac{2\pi a^{2}}{e}
			      \operatorname{artanh}e.
		      \]

		      因此，当 $a<b$ 时，
		      \[
			      \boxed{
				      S
				      =
				      2\pi b^{2}
				      +
				      \frac{2\pi a^{2}}{e}
				      \operatorname{artanh}e
			      },
			      \qquad
			      e=\sqrt{1-\frac{a^{2}}{b^{2}}}.
		      \]

	\end{description}

	\medskip

	\textbf{\textcolor{CProof}{情形三：$a=b$}}

	\begin{description}[style=nextline, leftmargin=2.8em, labelsep=0.5em, itemsep=0.55em]

		\item[\textbf{步骤八（退化为球）：}]
		      当 $a=b=R$ 时，椭圆退化为圆，绕 $x$ 轴旋转得到半径为 $R$ 的球。

		      此时
		      \[
			      I
			      =
			      \int_{-R}^{R}R^{2}\,dx
			      =
			      2R^{3}.
		      \]

		      代入
		      \[
			      S=\frac{2\pi R}{R^{2}}I
		      \]
		      得到
		      \[
			      S
			      =
			      \frac{2\pi R}{R^{2}}\cdot 2R^{3}
			      =
			      4\pi R^{2}.
		      \]

		      因此，当 $a=b=R$ 时，
		      \[
			      \boxed{
				      S=4\pi R^{2}
			      }.
		      \]

	\end{description}

	\medskip

	\textbf{\textcolor{CProof}{结论回扣}}

	综上所述，椭圆
	\[
		\frac{x^{2}}{a^{2}}+\frac{y^{2}}{b^{2}}=1
	\]
	绕 $x$ 轴旋转形成的旋转椭球表面积为：

	\[
		\begin{cases}
			\displaystyle
			S
			=
			2\pi b^{2}
			+
			\frac{2\pi ab}{e}\arcsin e,
			 &
			\displaystyle
			a>b,\quad
			e=\sqrt{1-\frac{b^{2}}{a^{2}}},
			\\[1.2em]
			\displaystyle
			S
			=
			2\pi b^{2}
			+
			\frac{2\pi a^{2}}{e}\operatorname{artanh}e,
			 &
			\displaystyle
			a<b,\quad
			e=\sqrt{1-\frac{a^{2}}{b^{2}}},
			\\[1.2em]
			\displaystyle
			S=4\pi R^{2},
			 &
			a=b=R.
		\end{cases}
	\]

\end{proofbox}